% Generated from locale/tr/content/proof-theory/natural-deduction/rules-proofs.tex; do not hand-edit.


\olfileid[tr]{pt}{nat}{rp}

\olsection{Kurallar ve \usetoken{P}{proof}}

Bazı tanımlar vererek ve bazı terimleri sabitleyerek başlıyoruz.

\Log{N1c} ve~\Log{N1i}'de $!A \lif !B$'den söz eden iki kuralı ele alın.
\[
\bottomAlignProof
\AxiomC{$\Discharge{!A}{x}$}
\shortDeduce
\DeduceC{$!B$}
\DischargeRule{\Intro{\lif}}{x}
\UnaryInfC{$!A \lif !B$}
\DisplayProof
\qquad
\bottomAlignProof
\AxiomC{$!A \lif !B$}
\AxiomC{$!A$}
\RightLabel{\Elim{\lif}}
\BinaryInfC{$!B$}
\DisplayProof
\]
\Elim{\lif} kuralı yeterince basittir. Çizginin üstündeki !!{formula}s
\emph{öncüllerdir}. Soldaki !!{formula}, \emph{ana} !!{formula}~$!A \lif
!B$'dir. Bütün !!{elimination} kurallarında her zaman en solda böyle bir öncül
bulunur. Buna çıkarımın \emph{ana} öncülü denir. Bu durumdaki gibi başka
öncüller de varsa bunlara \emph{yan} öncüller denir. !!{introduction}
kurallarının öncüllerine de ana öncül denir. Çizginin altındaki !!{formula},
\emph{sonuçtur}.

\Intro{\lif} kuralının yalnızca bir öncülü vardır ve burada sonuç ana
formül~$!A \lif !B$'dir. Bütün !!{introduction} kuralları böyledir: sonuç ana
formüldür; onun ana işlemcisi ``ortaya sokulan'' işlemcidir.

$\Discharge{!A}{x}$ gösterimi şu anlama gelir. Bir !!a{proof} içinde
\Intro{\lif} çıkarımı, öncül olarak bir $!B$ !!a{formula}{}üne uygulanır.
$!B$ ile biten alt-!!{proof}{}ın başlangıcında birtakım !!{formula}s bulunur.
Bunlara \emph{varsayımlar} denir. $!A$ biçimindeki bir ya da daha çok
varsayımdan $!B$'nin bir !!^a{proof}{}ı, $!A$'nın $!B$'yi gerektirdiğinin bir
!!a{proof}{}ıdır. \Intro{\lif} kuralını uyguladığımızda $!A \lif !B$ sonucuna
varırız ve sonuç artık $!A$ varsayımına dayanmaz. Bunu göstermek için,
\Intro{\lif} çıkarımını içeren !!{proof} içindeki $!A$ biçimindeki varsayımlar
\emph{boşaltılır} ya da \emph{kapatılır}. Kural, hangi varsayımların
boşaltılabileceğini belirtir. Sonucu $!A \lif !B$ olan \Intro{\lif} durumunda
bunlar $!A$ biçimindeki, yani koşullunun ön bileşeniyle örtüşen varsayımlardır.
Hangi varsayımların hangi kuralda boşaltıldığını belirtmek bazen önemlidir;
bu, boşaltma etiketi~$x$ ile yapılır. Önemli olarak, varsayımları boşaltan
kurallar gerekli biçimdeki varsayımların \emph{tümünü}, hatta herhangi birini
boşaltmak \emph{zorunda değildir}. Örneğin birinci \Intro{\lif} hiçbir
varsayımı boşaltmadığı hâlde aşağıdaki doğrudur:
\[
\AxiomC{$\Discharge{!A}{y}$}
\DischargeRule{\Intro{\lif}}{x}
\UnaryInfC{$!B \lif !A$}
\DischargeRule{\Intro{\lif}}{y}
\UnaryInfC{$!A \lif (!B \lif !A)$}
\DisplayProof
\]
Bir kural hiçbir varsayımı boşaltmadığında boşaltma etiketini yazmayız
(burada yaptığımız gibi).

Şu !!{proof} içinde
\[
\AxiomC{$\Discharge{!A}{x}$}
\AxiomC{$\Discharge{!A}{y}$}
\RightLabel{\Intro{\land}}
\BinaryInfC{$!A \land !A$}
\RightLabel{\Elim{\land}}
\UnaryInfC{$!A$}
\DischargeRule{\Intro{\lif}}{x}
\UnaryInfC{$!A \lif !A$}
\DischargeRule{\Intro{\lif}}{y}
\UnaryInfC{$!A \lif (!A \lif !A)$}
\DisplayProof
\]
birinci \Intro{\lif} çıkarımı soldaki $!A^x$ varsayımını, ikincisi sağdaki
$!A^y$ varsayımını boşaltır. Başka bir deyişle $x$ ile etiketlenmiş bütün
varsayımlar $x$ etiketli \Intro{\lif} tarafından, $y$ ile etiketlenmiş bütün
varsayımlar da $y$ etiketli olan tarafından boşaltılır. Bunun işlemesi için
aynı etikete sahip bütün varsayımların aynı zamanda aynı !!{formula} olması
önemlidir.

Niceleyici kuralları biraz daha karmaşıktır. Örneğin \Log{N1i}'de $\lexists$
için kurallar şunlardır:
\[
\bottomAlignProof
\AxiomC{$!A(t)$}
\RightLabel{\Intro{\lexists}}
\UnaryInfC{$\lexists[x][!A(x)]$}
\DisplayProof
\qquad
\bottomAlignProof
\AxiomC{$\lexists[x][!A(x)]$} 
\AxiomC{$\Discharge{!A(c)}{x}$}
\shortDeduce
\DeduceC{$!B$}
\DischargeRule{\Elim{\lexists}}{x}
\BinaryInfC{$!B$} \DisplayProof 
\bottomAlignProof
\] \Intro{\lexists} kuralında öncül $!A(t)$'dir; burada $t$ kapalı bir terim,
yani değişken içermeyen bir terimdir. Bir çıkarım, bir $!A$ !!{formula}{}ü,
bir $x$ değişkeni ve kapalı bir $t$ terimi bulunup sonuç
$\lexists[x][!A]$, öncül de~$\Subst{!A}{t}{x}$ olduğunda doğrudur---yani
\Intro{\lexists} şematik kuralına uyar. \Elim{\lexists}, $!A(c)$ biçimindeki
varsayımları boşaltmamıza izin verir; yani her çıkarım için bir !!{constant}~$c$
vardır ve $\Subst{!A}{c}{x}$ biçimindeki varsayımlar boşaltılabilir. Tek bir
çıkarımda boşaltılan varsayımların tümü aynı~$c$'yi kullanmalıdır. Bu sabite
\Elim{\lexists} çıkarımının \emph{özdeğişkeni} denir. Çıkarım ancak $c$,
boşaltılan $!A(c)$ dışındaki bir açık varsayımda, $!B$ içinde ve $!A$ içinde
geçmiyorsa doğrudur. Buna \emph{özdeğişken koşulu} denir.

\Log{N1c} ve~\Log{N1i}'deki !!^{proof}s, çıkarım kuralları kullanılarak
aksiyomlardan tümevarımsal biçimde üretilen sonlu formül ağaçlarıdır. Her
yaprak, bir boşaltma etiketiyle etiketlenmiş olabilen bir varsayımdır; diğer
her !!{formula}, sistemin çıkarım kurallarından biri aracılığıyla hemen
üstündeki !!{formula}s{}den çıkar. Bir çıkarımın üstündeki ağaçta varsayım olarak
geçen bir !!{formula}, onun üstündeki bir çıkarım tarafından boşaltılmamışsa
(o çıkarımda) \emph{açık} denir.

\begin{defn}\ollabel{defn:proof}
\Log{N1c} ve~\Log{N1i}'deki \emph{!!^{proof}s}, tümevarımsal olarak aşağıdaki
biçimde oluşturulan sonlu !!{formula} ağaçlarıdır:
\begin{enumerate}
  \item\ollabel{defn:prf:assumption} $x$ ya da $y$ gibi bir boşaltma etiketiyle
  etiketlenmiş her !!{formula}, tek başına bir !!a{proof}{}tır. Bir
  !!a{proof}, tek bir !!{formula}{}den oluşuyorsa bu !!{formula} açık varsayım
  sayılır.
  \item\ollabel{defn:prf:one-prem} $R$, bir, iki ya da üç öncülü $!A_1$,
  $!A_2$, $!A_3$, sonucu~$!A$ olan bir kural ve $\delta_1$, $\delta_2$,
  $\delta_3$ sırasıyla $!A_1$, $!A_2$, $!A_3$ ile biten !!{proof}s ise
  sırasıyla
  \[
  \AxiomC{}
  \RightLabel{$\delta_1$}
  \DeduceC{$!A_1$}
  \RightLabel{$R$}
  \UnaryInfC{$!A$}
  \DisplayProof
\qquad
  \AxiomC{}
  \RightLabel{$\delta_1$}
  \DeduceC{$!A_1$}
  \AxiomC{}
  \RightLabel{$\delta_2$}
  \DeduceC{$!A_2$}
  \RightLabel{$R$}
  \BinaryInfC{$!A$}
  \DisplayProof
\qquad
  \AxiomC{}
  \RightLabel{$\delta_1$}
  \DeduceC{$!A_1$}
  \AxiomC{}
  \RightLabel{$\delta_2$}
  \DeduceC{$!A_2$}
  \AxiomC{}
  \RightLabel{$\delta_3$}
  \DeduceC{$!A_3$}
  \RightLabel{$R$}
  \TrinaryInfC{$!A$}
  \DisplayProof
  \]
  birer !!a{proof}{}tır; yeter ki bu !!{proof}s içindeki varsayım etiketleri
  uyumlu olsun (iki farklı !!{formula} aynı boşaltma etiketine sahip olmasın)
  ve çıkarımlar $R$'nin doğru uygulamaları olsun.
  \item En üst !!{formula} geçişleri, yeni !!{proof} içinde $\delta_1$,
  $\delta_2$ ya da~$\delta_3$'ün açık varsayımlarıysa, kural tarafından
  boşaltılan geçişler dışında, bu !!{proof}{}ın \emph{açık varsayımları} sayılır.
  Kural $x$ (ya da $x\, y$) ile etiketlenmişse \Intro{\lif} ve~\Intro{\lnot}
  için bunlar $\delta_1$ içinde $x$ ile etiketlenmiş bütün açık varsayımlardır;
  \Elim{\lexists} için $\delta_2$ içinde $x$ ile etiketlenmiş bütün açık
  varsayımlardır; \Elim{\lor} için $\delta_2$ içinde $x$ ile ve $\delta_3$
  içinde $y$ ile etiketlenmiş bütün açık varsayımlardır.
\end{enumerate}
\end{defn}

Tümevarımsal kuruluşta kullanılan !!{proof}s{}a, ilgili !!a{proof}{}ın
\emph{alt-!!{proof}s{}ı} denir. Bir çıkarımın öncülleriyle biten alt-!!{proof}s{}a
çıkarımın \emph{doğrudan alt-!!{proof}s{}ı} denir. \Log{N1c} ve \Log{N1i}
kuralları \olref{tab:N1}'de verilmiştir.

\allmodulessubfile{OLP-0666}

\begin{defn}
\emph{!!{height}}~$\pheight{\delta}$ değeri, !!a{proof}~$\delta$ için
tümevarımsal olarak şöyle tanımlanır.
\begin{enumerate}
  \item $\delta$ bir varsayımdan oluşuyorsa $\pheight{\delta} = 0$.
  \item $\delta$, doğrudan alt-!!{proof}{}ı~$\delta_1$ olan tek öncüllü bir
  çıkarımla bitiyorsa $\pheight{\delta} = \pheight{\delta_1} + 1$.
  \item $\delta$, doğrudan alt-!!{proof}s{}ı~$\delta_1$ ve~$\delta_2$ olan iki
  öncüllü bir çıkarımla bitiyorsa $\pheight{\delta} =
  \max(\pheight{\delta_1}, \pheight{\delta_2}) + 1$.
  \item $\delta$, doğrudan alt-!!{proof}s{}ı~$\delta_1$, $\delta_2$ ve
  $\delta_3$ olan \Elim{\lor} ile bitiyorsa $\pheight{\delta} =
  \max(\pheight{\delta_1}, \pheight{\delta_2}, \pheight{\delta_3}) + 1$.
\end{enumerate}
Bir $S$ sekantının yüksekliği $\le n$ olan bir !!a{proof}{}ı varsa
$\Proves[n] S$ yazarız.
\end{defn}

\begin{ex}
\Log{N1i}'de herhangi bir !!{formula}, açık varsayım olarak kendisinden elde
edilen bir !!a{proof} sayılır. Öyleyse
\[
!A \land !B^x
\]
bir $\delta_1$ !!a{proof}{}ıdır. Yüksekliği~$0$'dır. $\delta_1$'den,
\Elim{\land}'ın iki sürümünden birini uygulayarak iki $\delta_2$ ve~$\delta_3$
!!{proof}{}ı elde edebiliriz:
\[
\AxiomC{$!A \land !B^x$}
\RightLabel{\Elim{\land}[2]}
\UnaryInfC{$!B$}
\DisplayProof
\qquad
\AxiomC{$!A \land !B^x$}
\RightLabel{\Elim{\land}[1]}
\UnaryInfC{$!A$}
\DisplayProof
\]
$\delta_2$ ve~$\delta_3$'teki son çıkarımın doğrudan alt-!!{proof}{}ı aynıdır:
$!A \land !B$. Her iki !!{proof} da $!A\land !B$ geçişini açık varsayım
olarak içerir. $\pheight{\delta_1} = 0$ ve
$\pheight{\delta_2} = \pheight{\delta_3} = 1$ olur.

\Intro{\land} kuralı, $!B \land !A$ biçimindeki bir sonucu gerekçelendirmek
için $!B$ ve $!A$ biçiminde iki öncül gerektirir. Dolayısıyla aşağıdaki bir
$\delta_4$ !!a{proof}{}ıdır:
\[
\AxiomC{$!A \land !B^x$}
\RightLabel{\Elim{\land}[2]}
\UnaryInfC{$!B$}
\LeftSubproofLabel{$\delta_2$}
\AxiomC{$!A \land !B^x$}
\RightLabel{\Elim{\land}[1]}
\UnaryInfC{$!A$}
\RightSubproofLabel{$\delta_3$}
\RightLabel{\Intro{\land}}
\BinaryInfC{$!B \land !A$}
\DisplayProof
\]
Yüksekliği $\pheight{\delta_4} = \max(\pheight{\delta_2},
\pheight{\delta_3})+1 = \max(1,1) + 1 = 2$'dir. Varsayım olarak iki
$!A \land !B$ geçişi vardır ve ikisi de açıktır.

Son olarak $\delta_5$'i elde etmek için \Intro{\lif} uygulayabiliriz:
\[
\AxiomC{$\Discharge{!A \land !B}{x}$}
\RightLabel{\Elim{\land}[2]}
\UnaryInfC{$!B$}
\LeftSubproofLabel{$\delta_2$}
\AxiomC{$\Discharge{!A \land !B}{x}$}
\RightLabel{\Elim{\land}[1]}
\UnaryInfC{$!A$}
\RightSubproofLabel{$\delta_3$}
\RightLabel{\Intro{\land}}
\BinaryInfC{$!B \land !A$}
\DischargeRule{\Intro{\lif}}{x}
\UnaryInfC{$(!A \land !B) \lif (!B \land !A)$}
\DisplayProof
\]
Yüksekliği $\pheight{\delta_4} + 1 = 3$'tür. \Intro{\lif} çıkarımı $x$ ile
etiketlenmiş $!A \land !B$ biçimindeki bütün varsayım geçişlerini, yani
doğrudan alt ispatın daha önce açık olan iki varsayımını da boşaltır. Bunlar
$\delta_5$'in tek varsayımları olduğundan hiçbir açık varsayımı kalmaz;
dolayısıyla bir !!a{proof}{}tır; ispat ettiği son-!!{formula}, $(!A \land !B)
\lif (!B \land !A)$'dır.
\end{ex}

